\documentclass[10pt]{article}

% ══════════════════════════════════════════════════════════════
%  PAQUETES BÁSICOS PARA ESTE TEST VISUAL
% ══════════════════════════════════════════════════════════════
\usepackage[utf8]{inputenc}
\usepackage[T1]{fontenc}
\usepackage[spanish]{babel}
\usepackage[margin=1.5cm, landscape]{geometry}
\usepackage{tikz}
\usetikzlibrary{matrix, positioning, shadows}
\usepackage{xcolor}
\usepackage{helvet}
\renewcommand{\familydefault}{\sfdefault}

% ══════════════════════════════════════════════════════════════
%  QUÍMICA (Para los ejemplos)
% ══════════════════════════════════════════════════════════════
\usepackage{chemfig}
\usepackage[version=4]{mhchem}

% ══════════════════════════════════════════════════════════════
%  DEFINICIÓN DE COLORES (PALETA SÓLIDA Y LEGIBLE)
% ══════════════════════════════════════════════════════════════

% --- Fundamentales (Texto normal) ---
\colorlet{elemH}{black}
\colorlet{elemC}{black}

% --- No Metales y Gases Nobles ---
\colorlet{elemN}{blue!85!black}
\colorlet{elemO}{red!85!black}
\colorlet{elemP}{orange!85!black}
\colorlet{elemS}{yellow!80!black}
\colorlet{elemSe}{orange!70!black}
\colorlet{elemTe}{orange!85!black}
\colorlet{elemPo}{brown!80!black}

\colorlet{elemF}{green!60!black}
\colorlet{elemCl}{green!80!black}
\colorlet{elemBr}{brown!80!black}
\colorlet{elemI}{purple!85!black}
\colorlet{elemAt}{purple!60!black}

\colorlet{elemHe}{cyan!80!black}
\colorlet{elemNe}{cyan!85!black}
\colorlet{elemAr}{teal!85!black}
\colorlet{elemKr}{teal!75!black}
\colorlet{elemXe}{teal!60!black}
\colorlet{elemRn}{teal!40!black}

% --- Metales Alcalinos (MolView: Violeta claro -> Morado -> Rosa Magenta) ---
\definecolor{elemLi}{RGB}{160, 120, 220}
\definecolor{elemNa}{RGB}{140, 80, 230}
\definecolor{elemK}{RGB}{170, 60, 240}
\definecolor{elemRb}{RGB}{210, 0, 200}
\definecolor{elemCs}{RGB}{210, 0, 140}
\definecolor{elemFr}{RGB}{240, 0, 180}

% --- Metales Alcalinotérreos (MolView: Verde olivo -> Verde vibrante) ---
\definecolor{elemBe}{RGB}{100, 120, 30}
\definecolor{elemMg}{RGB}{60, 150, 50}
\definecolor{elemCa}{RGB}{50, 160, 50}
\definecolor{elemSr}{RGB}{40, 170, 40}
\definecolor{elemBa}{RGB}{30, 130, 30}
\definecolor{elemRa}{RGB}{20, 110, 20}

% --- Grupo 13 (Cafés/Rosados) ---
\colorlet{elemB}{brown!60!black}
\colorlet{elemAl}{brown!70!black}
\colorlet{elemGa}{brown!80!black}
\colorlet{elemIn}{brown!90!black}
\colorlet{elemTl}{red!50!black}

% --- Grupo 14 (Grises Oscuros) ---
\colorlet{elemSi}{gray!80!black}
\colorlet{elemGe}{gray!70!black}
\colorlet{elemSn}{gray!60!black}
\colorlet{elemPb}{gray!50!black}

% --- Grupo 15 (Azules/Violetas) ---
\colorlet{elemAs}{violet!70!black}
\colorlet{elemSb}{purple!80!black}
\colorlet{elemBi}{purple!90!black}

% --- Metales de Transición (Tonos metálicos oscuros) ---
\colorlet{elemSc}{blue!50!black}
\colorlet{elemTi}{blue!60!black}
\colorlet{elemV}{blue!70!black}
\colorlet{elemCr}{blue!80!black}
\colorlet{elemMn}{blue!90!black}
\colorlet{elemFe}{orange!60!black}
\colorlet{elemCo}{blue!70!black}
\colorlet{elemNi}{green!70!black}
\colorlet{elemCu}{orange!80!black}
\colorlet{elemZn}{teal!70!black}

\colorlet{elemY}{blue!50!black}
\colorlet{elemZr}{blue!60!black}
\colorlet{elemNb}{blue!70!black}
\colorlet{elemMo}{blue!80!black}
\colorlet{elemTc}{blue!90!black}
\colorlet{elemRu}{orange!60!black}
\colorlet{elemRh}{blue!70!black}
\colorlet{elemPd}{green!70!black}
\colorlet{elemAg}{gray!70!black}
\colorlet{elemCd}{teal!60!black}

\colorlet{elemHf}{blue!50!black}
\colorlet{elemTa}{blue!60!black}
\colorlet{elemW}{blue!70!black}
\colorlet{elemRe}{blue!80!black}
\colorlet{elemOs}{orange!60!black}
\colorlet{elemIr}{blue!70!black}
\colorlet{elemPt}{green!70!black}
\colorlet{elemAu}{yellow!80!black}
\colorlet{elemHg}{teal!60!black}

% --- Lantánidos (MolView: Azul cielo -> Grisáceo -> Verde -> Teal) ---
\definecolor{elemLa}{RGB}{80, 160, 255}
\definecolor{elemCe}{RGB}{150, 180, 180}
\definecolor{elemPr}{RGB}{80, 180, 100}
\definecolor{elemNd}{RGB}{80, 180, 100}
\definecolor{elemPm}{RGB}{80, 180, 100}
\definecolor{elemSm}{RGB}{80, 180, 100}
\definecolor{elemEu}{RGB}{80, 180, 100}
\definecolor{elemGd}{RGB}{80, 150, 150}
\definecolor{elemTb}{RGB}{80, 150, 150}
\definecolor{elemDy}{RGB}{80, 150, 150}
\definecolor{elemHo}{RGB}{80, 160, 120}
\definecolor{elemEr}{RGB}{80, 180, 100}
\definecolor{elemTm}{RGB}{80, 180, 100}
\definecolor{elemYb}{RGB}{80, 180, 100}
\definecolor{elemLu}{RGB}{80, 180, 100}

% --- Actínidos (MolView: Azul bebé -> Morado azulado -> Rosa) ---
\definecolor{elemAc}{RGB}{80, 170, 255}
\definecolor{elemTh}{RGB}{80, 170, 255}
\definecolor{elemPa}{RGB}{80, 170, 255}
\definecolor{elemU}{RGB}{80, 150, 255}
\definecolor{elemNp}{RGB}{80, 150, 255}
\definecolor{elemPu}{RGB}{80, 130, 255}
\definecolor{elemAm}{RGB}{100, 100, 255}
\definecolor{elemCm}{RGB}{130, 80, 255}
\definecolor{elemBk}{RGB}{160, 80, 255}
\definecolor{elemCf}{RGB}{255, 80, 200}
\definecolor{elemEs}{RGB}{255, 80, 200}
\definecolor{elemFm}{RGB}{255, 80, 200}
\definecolor{elemMd}{RGB}{255, 80, 200}
\definecolor{elemNo}{RGB}{255, 80, 200}
\definecolor{elemLr}{RGB}{255, 80, 200}

% --- Elementos Súper Pesados (Gris Oscuro Opaco) ---
\colorlet{elemRf}{gray!60!black}
\colorlet{elemDb}{gray!60!black}
\colorlet{elemSg}{gray!60!black}
\colorlet{elemBh}{gray!60!black}
\colorlet{elemHs}{gray!60!black}
\colorlet{elemMt}{gray!60!black}
\colorlet{elemDs}{gray!60!black}
\colorlet{elemRg}{gray!60!black}
\colorlet{elemCn}{gray!60!black}
\colorlet{elemNh}{gray!60!black}
\colorlet{elemFl}{gray!60!black}
\colorlet{elemMc}{gray!60!black}
\colorlet{elemLv}{gray!60!black}
\colorlet{elemTs}{gray!60!black}
\colorlet{elemOg}{gray!60!black}

% ══════════════════════════════════════════════════════════════
%  MACROS PARA COLOREAR (Usando \textnormal para mantener la fuente)
% ══════════════════════════════════════════════════════════════
% Para texto normal: \el{O}
\newcommand{\el}[1]{\textcolor{elem#1}{\textnormal{#1}}}

% Para usar dentro de \ce{} de mhchem: \cel{O}
\newcommand{\cel}[1]{\textcolor{elem#1}{\textnormal{#1}}}

% Para usar dentro de \chemfig{}: \cfel{O}
\newcommand{\cfel}[1]{\textcolor{elem#1}{\textnormal{#1}}}

% Fix para babel español con chemfig
\AtBeginDocument{\shorthandoff{>}}

% ══════════════════════════════════════════════════════════════
%  DOCUMENTO
% ══════════════════════════════════════════════════════════════
\begin{document}

\begin{center}
    {\Huge \textbf{Test Visual de Colores: Tabla Periódica}}\\[0.5cm]
    {\large Paleta CPK sólida y gradientes por familia (Estilo MolView)}
\end{center}

\vspace{0.5cm}

\begin{center}
\begin{tikzpicture}[
    node distance=1.5mm,
    cell/.style={
        rectangle, 
        rounded corners=2pt, 
        minimum width=1.2cm, 
        minimum height=0.9cm, 
        font=\bfseries\large, 
        align=center,
        draw=black!30,
        thick
    }
]

% Fila 1
\matrix (p1) [matrix of nodes, nodes={cell}, column sep=2mm, row sep=2mm] {
    \el{H} & & & & & & & & & & & & & & & & & \el{He} \\
};

% Fila 2
\matrix (p2) [below=of p1, matrix of nodes, nodes={cell}, column sep=2mm] {
    \el{Li} & \el{Be} & & & & & & & & & & & \el{B} & \el{C} & \el{N} & \el{O} & \el{F} & \el{Ne} \\
};

% Fila 3
\matrix (p3) [below=of p2, matrix of nodes, nodes={cell}, column sep=2mm] {
    \el{Na} & \el{Mg} & & & & & & & & & & & \el{Al} & \el{Si} & \el{P} & \el{S} & \el{Cl} & \el{Ar} \\
};

% Fila 4
\matrix (p4) [below=of p3, matrix of nodes, nodes={cell}, column sep=2mm] {
    \el{K} & \el{Ca} & \el{Sc} & \el{Ti} & \el{V} & \el{Cr} & \el{Mn} & \el{Fe} & \el{Co} & \el{Ni} & \el{Cu} & \el{Zn} & \el{Ga} & \el{Ge} & \el{As} & \el{Se} & \el{Br} & \el{Kr} \\
};

% Fila 5
\matrix (p5) [below=of p4, matrix of nodes, nodes={cell}, column sep=2mm] {
    \el{Rb} & \el{Sr} & \el{Y} & \el{Zr} & \el{Nb} & \el{Mo} & \el{Tc} & \el{Ru} & \el{Rh} & \el{Pd} & \el{Ag} & \el{Cd} & \el{In} & \el{Sn} & \el{Sb} & \el{Te} & \el{I} & \el{Xe} \\
};

% Fila 6
\matrix (p6) [below=of p5, matrix of nodes, nodes={cell}, column sep=2mm] {
    \el{Cs} & \el{Ba} & \el{La} & \el{Hf} & \el{Ta} & \el{W} & \el{Re} & \el{Os} & \el{Ir} & \el{Pt} & \el{Au} & \el{Hg} & \el{Tl} & \el{Pb} & \el{Bi} & \el{Po} & \el{At} & \el{Rn} \\
};

% Fila 7
\matrix (p7) [below=of p6, matrix of nodes, nodes={cell}, column sep=2mm] {
    \el{Fr} & \el{Ra} & \el{Ac} & \el{Rf} & \el{Db} & \el{Sg} & \el{Bh} & \el{Hs} & \el{Mt} & \el{Ds} & \el{Rg} & \el{Cn} & \el{Nh} & \el{Fl} & \el{Mc} & \el{Lv} & \el{Ts} & \el{Og} \\
};

% Espacio para Lantánidos y Actínidos
\node[below=1.5cm of p7] (space) {};

% Lantánidos
\matrix (lan) [below=of space, matrix of nodes, nodes={cell}, column sep=2mm] {
    \el{La} & \el{Ce} & \el{Pr} & \el{Nd} & \el{Pm} & \el{Sm} & \el{Eu} & \el{Gd} & \el{Tb} & \el{Dy} & \el{Ho} & \el{Er} & \el{Tm} & \el{Yb} & \el{Lu} \\
};

% Actínidos
\matrix (act) [below=of lan, matrix of nodes, nodes={cell}, column sep=2mm] {
    \el{Ac} & \el{Th} & \el{Pa} & \el{U} & \el{Np} & \el{Pu} & \el{Am} & \el{Cm} & \el{Bk} & \el{Cf} & \el{Es} & \el{Fm} & \el{Md} & \el{No} & \el{Lr} \\
};

% Etiquetas laterales
\node[left=0.5cm of lan, font=\bfseries\large, text=gray] {Lantánidos};
\node[left=0.5cm of act, font=\bfseries\large, text=gray] {Actínidos};

\end{tikzpicture}
\end{center}

\vspace{1cm}

\begin{center}
\begin{tabular}{ll}
    \textbf{Notas de la paleta:} & \\
    H y C: & Negro (Texto normal) \\
    Alcalinos (Li-Fr): & Violeta claro $\to$ Morado $\to$ Rosa Magenta \\
    Alcalinotérreos (Be-Ra): & Verde olivo $\to$ Verde vibrante \\
    Metales de Transición: & Azules y tierras oscuras \\
    Gases Nobles (He-Rn): & Cian profundo $\to$ Verde azulado (teal) oscuro \\
    Lantánidos: & Azul cielo $\to$ Gris verdoso $\to$ Verde $\to$ Teal \\
    Actínidos: & Azul bebé $\to$ Azul profundo $\to$ Morado $\to$ Rosa \\
\end{tabular}
\end{center}

\vspace{1.5cm}

% ══════════════════════════════════════════════════════════════
%  EJEMPLOS DE USO
% ══════════════════════════════════════════════════════════════

\begin{center}
    {\Large \textbf{Ejemplos de uso en \texttt{mhchem} y \texttt{chemfig}}}
\end{center}

\vspace{0.5cm}

\noindent
\textbf{1. En \texttt{mhchem} (\texttt{\textbackslash ce\{\}}):}
Usa el macro \texttt{\textbackslash cel\{\}} para colorear cada elemento.
\begin{itemize}
    \item Ácido sulfúrico: \ce{\cel{H}2\cel{S}\cel{O}4}
    \item Permanganato de potasio: \ce{\cel{K}\cel{Mn}\cel{O}4}
    \item Etanol: \ce{\cel{C}\cel{H}3-\cel{C}\cel{H}2-\cel{O}\cel{H}}
    \item Reacción: \ce{\cel{Na}\cel{Cl} + \cel{Ag}\cel{N}\cel{O}3 -> \cel{Ag}\cel{Cl} v + \cel{Na}\cel{N}\cel{O}3}
\end{itemize}

\vspace{0.5cm}

\noindent
\textbf{2. En \texttt{chemfig} (\texttt{\textbackslash chemfig\{\}}):}
Usa el macro \texttt{\textbackslash cfel\{\}} o \texttt{\textbackslash textcolor\{elemX\}\{\}} directamente dentro del nodo.
\begin{center}
\begin{tabular}{ccc}
    \textbf{Etanol} & \textbf{Acetona} & \textbf{Fenol} \\
    \vspace{0.2cm}
    \chemfig{\cfel{C}H_3-\cfel{C}H_2-\cfel{O}H} &
    \chemfig{\cfel{C}H_3-\cfel{C}(=[2]\cfel{O})-\cfel{C}H_3} &
    \chemfig{*6(=-=-(-\cfel{O}H)=-)} \\
\end{tabular}
\end{center}

\vspace{0.5cm}

\noindent
\textbf{3. Estructuras más complejas:}
\begin{center}
\begin{tabular}{cc}
    \textbf{Glicina} & \textbf{Cafeína} \\
    \vspace{0.2cm}
    \chemfig{\cfel{N}H_2-\cfel{C}H_2-\cfel{C}(=[2]\cfel{O})-\cfel{O}H} &
    \chemfig{*5(=-=-(-[::+0]\cfel{N}(-[::+60]\cfel{C}H_3)-[::-60]\cfel{C}(=[::+0]\cfel{O})-[::+60]\cfel{N}(-[::+60]\cfel{C}H_3))=-)} \\
\end{tabular}
\end{center}

\vspace{1.5cm}

% ══════════════════════════════════════════════════════════════
%  COMPARACIÓN VISUAL 
% ══════════════════════════════════════════════════════════════
\begin{center}
    {\Large \textbf{Comparación Tipográfica Final}}\\[0.5cm]
    
    {\Huge
    % Fórmula 1: Entorno normal
    \ce{H2O}
    
    \vspace{0.5cm}
    
    % Fórmula 2: Con nuestros macros de color
    \ce{\cel{H}2\cel{O}}
    }

\end{center}
\end{document}